\documentclass[11pt]{article}
\usepackage[margin=1in]{geometry}
\usepackage{amsmath,amssymb,amsthm,booktabs}
\usepackage[T1]{fontenc}
\usepackage{lmodern}
\usepackage{microtype}
\usepackage{hyperref}
\hypersetup{colorlinks=true,linkcolor=blue,urlcolor=blue}
\setlength{\emergencystretch}{3em}
\newtheorem{theorem}{Theorem}
\newtheorem{proposition}{Proposition}
\newcommand{\cI}{\mathcal I}
\newcommand{\norm}[1]{\left\lVert #1\right\rVert}
\title{A Growing-Cutoff and Convex-Profile Transfer Audit for the V59 Residual}
\author{Liang Wang\\School of Mathematics and Statistics\\Huazhong University of Science and Technology (HUST)\\Wuhan, China}
\date{26 August 2026}
\begin{document}
\maketitle
\begin{abstract}
TPC-268 showed that a fixed finite comparison cutoff can move a literal V59
residual across the quarter threshold. We now replace that fixed choice by the
registered finite proxy $z_N=\lfloor\log N\rfloor$ and move the kernel along
the convex family $K_\theta=(1-\theta)K_{H,1}+\theta K_{H,2}$ generated by
two normalized nonnegative profiles. The physical prime shell, masks, deleted
diagonal, source coefficient, and rank-three projection remain unchanged.
Outward rational intervals classify twelve rows: eight contractions and four
obstructions. At the same growing-cutoff clock, the central profile path has
$\theta=9/10$ above and $\theta=24/25$ below $1/4$. This is a finite
model-relative transfer obstruction, not an asymptotic V59 estimate.
\end{abstract}
\section{Position and claim firewall}
The residual object left by the typed TPC-266 chain is a literal prime-shell
quantity whose radius and phase must be controlled at growing scale. TPC-267
instantiated that object at twelve finite rows with a $z=2$ comparison and
found a strict quarter contraction. TPC-268 then demonstrated that changing
declared finite cutoff, clock, or kernel parameters can reverse the verdict.
The present paper asks the next natural question: does a declared scale-following
cutoff together with a convex profile path remove that finite sensitivity?

We use the following status label:
\begin{center}
\texttt{NUMERICALLY\_CERTIFIED\_FINITE\_GROWING\_CUTOFF\_PROFILE\_TRANSFER}
\end{center}
The word ``growing'' below refers to the finite proxy registry unless it is
explicitly qualified as source-level or asymptotic. No endpoint exponent is
inferred from the table.
\section{The finite source and operator}
Let $N$ be even and set
\[
 \cI_N=\{N/2+1,\ldots,N\},\qquad
 \mathcal Q_Q=\{q:q\text{ prime},\ Q<q\leq2Q\}.
\]
For an integer $z\geq2$, let $b_N^{(z)}$ be the finite shifted-prime
comparison from TPC-268:
\[
 b_N^{(z)}(u)=C_2^{(z)}1_{2\nmid u}
 \prod_{p\leq z}\frac p{p-1}
 \prod_{\substack{p\mid u\\p>z}}\frac{p-1}{p-2},
\]
where the value is zero when $u+2$ has a prime factor at most $z$, and
$C_2^{(z)}=\prod_{p>z}(1-(p-1)^{-2})$. Define
\[
 \beta_N(t)=\frac{\Lambda(t)}{\log t}
 -\sum_{\substack{d\mid t\\d^{400}\leq N^{133}}}\mu(d),
 \qquad w_N(u)=\Lambda(u+2)-b_N^{(z_N)}(u).
\]
The finite cutoff proxy is
\[
 z_N=\lfloor\log N\rfloor,
 \qquad (z_{64},z_{96},z_{128},z_{192},z_{256},z_{384})=(4,4,4,5,5,5).
\]
For $s=1,2$, use $K_{H,s}(h)=(1+(h/H)^2)^{-s}$. These are the two
finite kernel representatives inherited from the previous audit. If $A_s$
denotes (with the same shell, masks, and deleted diagonal)
\[
 A_s(u,t)=1_{u\ne t}\sum_{q\in\mathcal Q_Q}qK_{H,s}(u-t)1_{q\nmid ut}
 \left(1_{u\equiv t\pmod q}-\frac1{q-1}\right),
\]
then for rational $0\leq\theta\leq1$ define
\[
 A_\theta=(1-\theta)A_1+\theta A_2,\qquad
 g_\theta=A_\theta\beta_N.
\]
\section{Exact profile transfer}
Partition $\cI_N$ into four equal consecutive blocks of size $B=N/8$ and
use the orthogonal contrasts
\[
 c_0=(1,1,-1,-1),\quad c_1=(1,-1,0,0),\quad
 c_2=(0,0,1,-1).
\]
If $W_j$ and $G_j$ are the block sums of $w_N$ and $g_\theta$, put
\[
 C_\theta=\langle w_N,g_\theta\rangle,\qquad
 C_{3,\theta}=\frac{W_0G_0}{4B}+\frac{W_1G_1}{2B}+\frac{W_2G_2}{2B},
 \qquad C_{\perp,\theta}=C_\theta-C_{3,\theta}.
\]
Let
\[
 R_\theta^2=\norm{(I-P_3)w_N}^2
                 \norm{(I-P_3)g_\theta}^2,\qquad
 \rho_\theta^2=\frac{|C_{\perp,\theta}|^2}{R_\theta^2}.
\]
\begin{proposition}[affine transfer identity]
For every finite row and rational $\theta$,
$g_\theta=(1-\theta)g_1+\theta g_2$ and
$C_{\perp,\theta}=(1-\theta)C_{\perp,1}+\theta C_{\perp,2}$.
The projection identity $C_\theta=C_{3,\theta}+C_{\perp,\theta}$ is exact.
\end{proposition}
\begin{proof}
All entries of $A_1,A_2$ are rational on a finite row. Linearity of matrix
multiplication and of the orthogonal projection gives the first two identities.
The three contrasts are pairwise orthogonal, so the displayed denominators are
their squared norms. Subtracting the projected cross-Gram proves the last
identity. The denominator $R_\theta^2$ is quadratic in $\theta$, so the
proposition deliberately makes no monotonicity assertion about $\rho_\theta$.
\end{proof}
\section{Interval certificate}
The Euler products are multiplied through $P=50000$. The positive tail bound
\[
 \prod_{p>P}\left(1-\frac1{(p-1)^2}\right)
 \geq1-\sum_{p>P}\frac1{(p-1)^2}>1-\frac1{P-1}
\]
gives rational endpoints. Logarithms use 100-digit decimal intervals with a
$10^{-25}$ guard, and all endpoints are rounded outward to a $10^{-30}$ grid.
The decision is made before taking a square root:
\[
 \rho^2_{\rm hi}<\frac1{16}\Rightarrow\text{contraction},\qquad
 \rho^2_{\rm lo}>\frac1{16}\Rightarrow\text{obstruction}.
\]
The producer computes the two exact rational operator outputs separately and
forms their rational convex combination. The independent checker uses an
independent floating-point implementation; it is an audit rather than a
substitute for the interval enclosure.
\begin{theorem}[finite growing-cutoff/profile transfer]
The certificate contains twelve threshold-separated rows, eight contractions
and four obstructions. Its six base rows use $\theta=0$ and the cutoff
$z_N=\lfloor\log N\rfloor$. In the central row $(N,H,Q)=(64,15,4)$,
\[
 \begin{array}{c|c|c}
 \theta&\rho^2\text{ interval}&\text{classification}\\ \hline
 9/10&[0.0634078324659,\,0.0634208686352]&\text{obstruction}\\
 24/25&[0.0622500850692,\,0.0622630874560]&\text{contraction}\\
 1&[0.0614513775060,\,0.0614643508731]&\text{contraction}
 \end{array}
\]
so the convex profile path crosses the quarter threshold while the cutoff,
clock, shell, and projection are fixed.
\end{theorem}
\begin{proof}
The finite identities are the proposition. The interval producer enumerates the
registered rows and propagates rational endpoints through the projection and
quotient. Every contraction upper endpoint is below $1/16$ and every obstruction
lower endpoint is above $1/16$. The independent replay reproduces all twelve
strict inequalities and the stress audit checks the cutoff registry and the
profile flip. These checks establish the stated finite theorem.
\end{proof}
\begin{table}[t]
\centering
\caption{Selected rows. Values are rounded square-root upper endpoints; the
decision uses the stored $\rho^2$ interval.}
\small
\begin{tabular}{rrrrrll}
\toprule
N&H&Q&$z_N$&$\theta$&displayed $\rho$&class\\
\midrule
64&15&4&4&0&0.2735358969&obstruct\\
96&20&5&4&0&0.2996937679&obstruct\\
128&24&5&4&0&0.1248321291&contract\\
192&32&6&5&0&0.0072691733&contract\\
256&38&6&5&0&0.0338322886&contract\\
384&50&7&5&0&0.0310909465&contract\\
64&15&4&4&9/10&0.2518350028&obstruct\\
64&15&4&4&24/25&0.2495257250&contract\\
96&20&5&4&1/2&0.2985106861&obstruct\\
128&24&5&4&1/2&0.1421214943&contract\\
256&38&6&5&1/2&0.0157231342&contract\\
\bottomrule
\end{tabular}
\end{table}
\section{Interpretation and limits}
The finite proxy does two useful things. It removes the arbitrary fixed $z=2$
from the experiment and makes the profile dependence explicit through an
affine path. It also gives a sharp warning: a favorable endpoint profile does
not imply a profile-uniform phase bound, because $\rho_\theta$ is a quotient
with a changing denominator.

Three promotions are prohibited. The schedule $z_N=\lfloor\log N\rfloor$ is
not a proof of the source-level V59 cutoff theorem. The twelve rows do not
control $R_\theta$ at growing scale and hence pay no fixed-power credit. Finally,
the finite profile path is not an arithmetic $L^2$ estimate and does not close
the collective prime-shell reassembly in Gate B.
\section{Conclusion}
TPC-269 supplies a finite transfer obstruction with a new source-compatible
proxy: a growing cutoff and a convex profile family on the same physical
operator. The central $9/10$ versus $24/25$ flip shows that even this
canonicalized finite interface is not uniformly inside the quarter sector.
The next minimal task is cross-scale normalization of the residual radius, with
the profile and cutoff held source-compatible. The actual V59 radius, phase,
arithmetic $L^2$, strict $1/400$ payment, and twin-prime conclusion remain open.
\begin{thebibliography}{9}
\bibitem{hardywright} G.~H.~Hardy and E.~M.~Wright,
\emph{An Introduction to the Theory of Numbers}, 6th ed., Oxford University
Press, 2008.
\bibitem{montgomeryvaughan} H.~L.~Montgomery and R.~C.~Vaughan,
\emph{Multiplicative Number Theory I: Classical Theory}, Cambridge University
Press, 2007.
\bibitem{tpc268} L.~Wang, ``A finite cutoff-sensitivity obstruction for the
V59 residual phase,'' TPC-268 project release, 2026.
\end{thebibliography}
\end{document}
